\section{Results}
\label{sec:results}

\subsection{Calibrated guidance improves paired restoration}
\label{sec:res-paired}

On the fifteen development participants, matched calibration reached a participant-level $\rrmse$ of 0.4310, against 0.5738 for the unguided control and 0.6511 for population calibration (Table~\ref{tab:paired}). The improvement over the unguided control, 0.1428 (95\% bootstrap interval 0.1061--0.1834), was positive for all fifteen participants and isolates what the explicit guide contributes beyond the shared network and calibration features. The improvement over population calibration, 0.2202 (0.0929--0.4025), was positive for fourteen of fifteen. Population calibration is in fact worse than raw EEG (0.6511 versus 0.6076), and the reason is instructive: an averaged propagation matrix applied to a person it does not fit subtracts activity that was never artifact. A population estimate is therefore useful as a shrinkage target, where it stabilizes an individual estimate, and harmful as a substitute, where it drives a subtraction of its own. This observation shaped the fallback of Section~\ref{sec:method-calibration}.

The ablation matrix locates the same conclusion cell by cell (Table~\ref{tab:matrix}; all cells share one episode bank and noise draw, so the matrix is read as contrasts against the proposed cell). With a matched guide, replacing the matched calibration features by their population values costs essentially nothing (+0.005, interval $-0.004$ to $+0.018$): the guide carries the payload, and the features are auxiliary. Without the guide, matched features recover only a small part of the gap. The one large cell is the cautionary one: guiding with the \emph{unshrunk} 120-s ridge estimate ($\lambda=1$) is worse than the shrunk guide by 0.347 (interval $+0.004$ to $+1.025$; worse for 13 of 15 participants) --- the raw two-minute estimate is occasionally catastrophically wrong, and empirical Bayes shrinkage is what makes a two-minute calibration deployable at all.

\begin{table}[t]
\caption{Guide $\times$ calibration-feature ablation (one episode bank; contrasts vs the proposed cell, positive = worse than proposed; 95\% bootstrap intervals).}
\label{tab:matrix}
\centering
\small
\begin{tabular}{@{}llrl@{}}
\toprule
Guide & Features & vs (matched, matched) & worse \\
\midrule
matched & matched & --- & --- \\
matched & population & $+0.005\ (-0.004, +0.018)$ & 7/15 \\
none & matched & $+0.142\ (+0.080, +0.207)$ & 13/15 \\
none & population & $+0.158\ (+0.093, +0.223)$ & 13/15 \\
population & population & $+0.186\ (+0.095, +0.298)$ & 14/15 \\
matched, unshrunk ($\lambda{=}1$) & matched & $+0.347\ (+0.004, +1.025)$ & 13/15 \\
\bottomrule
\end{tabular}
\end{table}

The fallback choice was then measured directly. The reliability criterion rejected the calibration in four development cells. On the subtraction pathway, replacing the population-substituted subtraction by no subtraction improved exactly those cells by 0.2537 (0.0070--0.5005) and the pooled error from 0.4340 to 0.4227; within the diffusion system the same change gave 0.0614 ($-$0.0258--0.1919) on the four cells --- a positive point estimate whose interval spans zero at $n=4$ because the learned corrections partially absorb a bad guide, a robustness margin the linear pathway does not have. On natural recordings from the rejected cells, withholding the guide raised retention from 0.917 to 0.984 while attenuation fell from 0.41 to 0.13 dB: the system stops subtracting where it cannot trust its calibration. Rerunning the full development comparison under this final rule left the headline effect unchanged, 0.1470 (0.1121--0.1871) with fifteen of fifteen positive; the two rules coincide on every cell whose calibration is accepted.

The eight held-out participants were then evaluated once, with frozen models, a frozen population registry, and the comparison fixed in advance. Matched calibration reached 0.5351 against 0.6888 for the unguided control, an improvement of 0.1537 (0.0725--0.2504) with seven of eight participants positive, close to the development effect despite the higher absolute error of this cohort. The comparison floor is itself demanding: on these participants the unguided control already improved on raw EEG by 0.086, so the calibrated gain sits on top of everything the population model and calibration features deliver. Figure~\ref{fig:denoise} shows the paired waveforms, the per-participant contrasts, and the scalp topography of the improvement, which is ocular where it should be (largest at Fp1, Fp2, and AFz; smallest at the ear electrodes).

\begin{table*}[t]
\caption{Participant-level paired restoration. Lower $\rrmse$ is better; improvements are participant-level mean differences with 95\% bootstrap intervals. The two cohorts are never pooled.}
\label{tab:paired}
\centering
\small
\begin{tabularx}{\textwidth}{@{}lYcYYc@{}}
\toprule
Cohort & Condition & $\rrmse$ & Comparison & Improvement (95\% interval) & Positive \\
\midrule
\multirow{4}{*}{Development ($n=15$)}
 & Raw EEG & 0.6076 & -- & -- & -- \\
 & Population calibration & 0.6511 & Matched vs population & 0.2202 (0.0929--0.4025) & 14/15 \\
 & Unguided control & 0.5738 & Matched vs unguided & 0.1428 (0.1061--0.1834) & 15/15 \\
 & Matched calibration & \textbf{0.4310} & -- & -- & -- \\
\midrule
\multirow{4}{*}{Held-out ($n=8$)}
 & Raw EEG & 0.7753 & -- & -- & -- \\
 & Population calibration & 0.7661 & -- & -- & -- \\
 & Unguided control & 0.6888 & Matched vs unguided & 0.1537 (0.0725--0.2504) & 7/8 \\
 & Matched calibration & \textbf{0.5351} & -- & -- & -- \\
\bottomrule
\end{tabularx}
\end{table*}

\begin{figure*}[t]
  \centering
  \figureplaceholder{56mm}{\textbf{Paired denoising and participant effects (design).} Panel A: a waveform strip for development participant sub-02 (ses-02, SSVEP cell, episode 1 --- the lowest-identifier participant, earliest clear ocular event) showing, on shared time and amplitude axes, the recorded EOG drive, the contaminated EEG, the paired reference, linear regression, the unguided control, and matched calibration for the three most anterior channels (Fp1, Fp2, AFz); arrays in \texttt{t6\_waveform\_exemplar\_dev.npz}. Panel B: paired per-participant $\rrmse$ lines from unguided to matched for all fifteen development participants. Panel C: the same for all eight held-out participants. Panel D: the scalp map of the per-channel improvement (\texttt{t6\_scalp\_improvement.npz}; Fp1 $+0.29$, Fp2 $+0.27$, AFz $+0.22$, ear electrodes near zero; the 14 ear channels are listed beside the 32-channel scalp map). All waveforms and points from stored evaluator outputs; nothing synthesized.}
  \caption{Planned figure: what calibrated guidance changes in the waveform, the per-participant effects behind Table~\ref{tab:paired}, and the ocular topography of the improvement.}
  \Description{Placeholder for waveform comparisons, paired participant plots, and an improvement scalp map.}
  \label{fig:denoise}
\end{figure*}

Against classical and deterministic references in the development protocol, the guided sampler matches the strongest calibrated point estimator --- matched diffusion 0.4310 versus calibrated linear regression 0.4353 --- with the one-step deterministic network at 0.4968; when both learned estimators are retrained with population calibration only, the deterministic network is ahead (0.5043 versus 0.5260). The calibrated operator sets the operating point; what the diffusion prior adds is everything around it: one frozen network that serves the guided, unguided, population, and fallback modes in-distribution, a measured robustness margin under a bad guide, and a conditional distribution over clean waveforms that Section~\ref{sec:res-uq} shows is better calibrated than a deterministic ensemble where distributions matter.

\subsection{Natural recordings, specificity, and on-line adaptation}
\label{sec:res-natural}

On natural development recordings, matched calibration attenuated 2.463 dB of EOG-coherent signal at a low-EOG retention of 0.843; population calibration attenuated 0.909 dB at retention 0.773, and the unguided control preserved the most (0.921) while attenuating the least (0.277 dB). The completed five-condition plane (Table~\ref{tab:natural}) adds the specificity conditions: a wrong person's matrix pushed through the gate behaves like population calibration (0.570 dB at 0.780), the same matrix without the gate actively injects signal (attenuation $-1.523$ dB at retention 0.334), and temporally shuffled EOG is nearly as destructive ($-0.384$ dB at 0.365). Matched calibration dominates population calibration on both axes at once, gaining 1.554 dB of attenuation (0.969--2.211; 14/15 participants), 0.0707 of retention (0.0382--0.1067; 13/15), and 0.0849 of EOG--EEG coherence reduction (0.0535--0.1211; 14/15). The paired-error gain is not obtained by subtracting more everywhere; it is obtained by subtracting the right thing.

The literature anchors make the same point from the other side. On this protocol ICA attenuates 2.98 dB at retention 0.704, ASR 2.81 dB at 0.623, and the calibrated eye-subspace projection 4.87 dB at 0.495 --- every anchor buys its attenuation below our 0.75 retention validity bar, removing real signal along with the artifact (the subspace projection also degrades the paired reconstruction to 0.827, worse than raw). These are context rows, not contests --- ICA sees the whole continuous record and needs no EOG; ASR failed its calibration on two cells, which its design cannot report --- but on the attenuation--retention plane the calibrated guided system is the only operating point with both axes healthy (Figure~\ref{fig:natural}).

\begin{table}[t]
\caption{Natural-recording operating points, development cohort. Natural recordings have no clean target; higher attenuation and higher retention are both preferred; the validity bar requires retention $\geq 0.75$.}
\label{tab:natural}
\centering
\small
\begin{tabular}{@{}lccc@{}}
\toprule
Condition & Attenuation (dB) & Retention & Coherence red. \\
\midrule
Unguided control & 0.277 & \textbf{0.921} & 0.021 \\
Population calibration & 0.909 & 0.773 & 0.105 \\
Matched calibration & 2.463 & 0.843 & 0.190 \\
Mismatched (gated) & 0.570 & 0.780 & 0.084 \\
Mismatched (ungated) & $-1.523$ & 0.334 & 0.029 \\
Shuffled EOG & $-0.384$ & 0.365 & 0.083 \\
\midrule
ICA (EOG-corr.) & 2.983 & 0.704 & 0.219 \\
ASR (120-s calib.) & 2.812 & 0.623 & 0.119 \\
Eye-subspace proj. & \textbf{4.869} & 0.495 & 0.213 \\
\bottomrule
\end{tabular}
\end{table}

The controls pin down what the calibration is and is not. Substituting another participant's matrix degrades reconstruction, and destroying the temporal alignment of the EOG removes most of the guide's value, so the benefit rides on the specific propagation of the specific person, delivered at the right time. The reliability criterion, however, is a check on estimation quality, not on ownership: a stable matrix from the wrong person passes it, and the accepted mismatched condition remains 0.0384 worse than population calibration (0.0089--0.0697). The system accordingly assumes that calibration and evaluation data come from the same person and session, which is how the acquisition protocol is designed.

Calibration also cannot be replaced by adapting during use. A recursive least-squares tracker that updates the propagation estimate causally during evaluation reached $\rrmse$ 0.944 after 10 s of adaptation, 0.654 after 30 s, 0.626 after 60 s, 0.601 after 120 s, and 1.051 after 240 s, never approaching the static calibrated value of 0.431; initializing the tracker wrongly made no lasting difference, since it converged to the same level within about 30 s. Evaluation-time data is dominated by exactly the artifacts being removed, which makes it poor material for estimating the artifact pathway. A two-minute dedicated calibration buys what continuous adaptation does not, which is why the system's adaptation loop runs through recalibration rather than through on-line updates. Figure~\ref{fig:natural} collects the natural operating points, the specificity controls, and the adaptation comparison.

\begin{figure*}[t]
  \centering
  \figureplaceholder{50mm}{\textbf{Natural behavior, specificity, and adaptation (design).} Panel A: operating points with low-EOG retention on the horizontal axis and EOG attenuation on the vertical axis for all nine conditions of Table~\ref{tab:natural} (diffusion-family conditions as filled markers, literature anchors as open markers, the 0.75 validity bar as a vertical line); matched should be the only point in the upper-right healthy region (arrays in \texttt{t5\_natural\_plane.npz} and \texttt{cpu\_reference\_rows.npz}). Panel B: per-participant matched-minus-population differences for attenuation, retention, and coherence reduction with bootstrap intervals. Panel C: $\rrmse$ of the recursive least-squares tracker against adaptation time (10, 30, 60, 120, 240 s) as a curve, with the static calibrated value as a horizontal line. All values from stored evaluator outputs.}
  \caption{Planned figure: matched calibration removes more while preserving more; classical anchors attenuate strongly but below the retention bar; the benefit is owner- and time-specific; on-line adaptation does not reach the calibrated operating point.}
  \Description{Placeholder for attenuation-retention operating points, specificity effects, and the adaptation-time curve.}
  \label{fig:natural}
\end{figure*}

\subsection{Calibration economics: the duration curve and the operator's lifetime}
\label{sec:res-loop}

How much calibration does the loop need, and how long does it live? Truncating the calibration prefix and recomputing everything closed-form gives the duration curve (Figure~\ref{fig:loop}). Sixty seconds already delivers a gain of 0.146 (0.091--0.210; 14/15) against 0.150 (0.098--0.209) at the full two minutes --- 97\% of the benefit at half the cost. Even a 30-second ridge estimate carries 0.124 (0.056--0.193) when the rejection floor is bypassed; with the rule active, 30 s is rejected on every cell by construction (the four 30-s sub-blocks that drive shrinkage and the reliability verdict require 60 s), and the system reading sits at the safe unguided point ($-0.016$, indistinguishable from zero) rather than degrading. At 60 s and beyond, the system reading is slightly \emph{above} the rule-off reading at every duration: the gate is not only a safety device, it nets a small positive by rejecting the few unreliable calibrations. The floor at 60 s is therefore exactly where the information budget allows it to be, and below it the system refuses instead of failing.

The lifetime has two faces. The operator itself drifts: re-estimated on successive 120-s windows of the same recording, the propagation matrix moves away from the calibration estimate by 86\% of its norm within two minutes and by 143\% at eight minutes, and the real-recording benefit follows --- the natural attenuation gain of matched over unguided decays from about 2.5 dB in the first evaluated windows to 1.5 dB (1.07--1.95) by the ten-minute mark. The guide's validity on \emph{content}, in contrast, persists: on fresh paired episodes drawn from early, mid, and late record thirds the gain is flat to rising (0.134, 0.143, 0.158). Together these give the adaptation loop its clock: what expires is the operator, not the mechanism, and the corrective action is a fresh two-minute calibration.

\begin{figure*}[t]
  \centering
  \figureplaceholder{46mm}{\textbf{Calibration economics (design).} Panel A: gain versus calibration duration (30/60/90/120 s) with bootstrap bands, two curves --- system reading (reliability rule active; the 30-s point sits at the unguided floor with a marker noting `gate fires 100\%') and rule-off reading (raw shrinkage curve); arrays in \texttt{t2\_duration\_curve.npz}. Panel B: relative operator displacement versus window start (0--480 s; \texttt{t4\_staleness.npz}) and, on a twin axis, the natural attenuation gain by window position (2.45/2.55/2.24/1.51 dB) --- the drift curve rises as the benefit falls. All values from stored evaluator outputs.}
  \caption{Planned figure: sixty seconds of calibration buys 97\% of the gain and the floor below it is a designed refusal, not a failure; the calibrated operator has a measurable lifetime that points the loop at recalibration.}
  \Description{Placeholder for the calibration-duration curve and the operator-lifetime curves.}
  \label{fig:loop}
\end{figure*}

\subsection{Predictive intervals: calibrated coverage, confirmed on held-out users}
\label{sec:res-uq}

The raw distribution of the 32 trajectories is too narrow: empirical 50/80/90\% coverages are 0.302, 0.479, and 0.558, with a continuous ranked probability score (CRPS) of 0.1516 (Table~\ref{tab:intervals}). A fold-held-out scalar temperature repairs coverage (0.625/0.802/0.851) at CRPS 0.1524. Adding the propagation-matrix uncertainty of Equation~\ref{eq:operator-sample} and then temperature-scaling reaches 0.622/0.802/0.853 at CRPS 0.1503, with the risk--coverage area improving from 0.0520 to 0.0504. The 80\% and 90\% levels are within 0.05 of nominal, and the calibrated intervals are not blunter than the raw ones, they are sharper: because the raw samples were underdispersed, widening them toward the truth lowers CRPS, and the physically motivated width does it better than temperature alone, with a lower CRPS (0.1503 versus 0.1524) and a smaller required temperature (2.30--2.55 versus 2.90--3.20). Across cells the operator term contributes a roughly constant share of the squared width (about 0.30); its value shows in the probabilistic score and the temperature it saves, not in a per-cell localization. The 50\% interval remains conservative (0.622), the one nominal level a single scalar does not align.

The held-out confirmation is the chapter's point. One scalar per policy was frozen on all fifteen development cells jointly (2.40 with the propagation term, 3.00 without) and applied unchanged to the eight held-out participants, evaluated once. Coverage held: 0.610/0.810/0.864 at the three nominal levels for the adopted policy, against 0.613/0.813/0.867 for temperature-only --- and the development ordering reproduced out of cohort, the propagation-augmented width again scoring better (CRPS 0.3800 versus 0.3824; absolute CRPS is larger than in development because held-out errors are larger, and is not comparable across cohorts) at the smaller temperature, with the raw samples again badly undercovered (0.460/0.543). Per-participant 80\% coverage spans 0.77--0.94 for six of eight participants, with sub-04 at 0.774 and one clear undercoverage at 0.456 for the cohort's most difficult participant. Every structural property of the interval mechanism --- undercoverage of raw samples, repair by one scalar, the advantage and the smaller temperature of the physically motivated width, the conservative 50\% level --- transferred to users the system had never seen, under settings chosen before their data was touched.

\begin{table}[t]
\caption{Predictive-interval quality ($K=32$). Development: temperatures fold-held-out. Held-out cohort: one scalar per policy frozen on development and applied unchanged, single pass. CRPS is cohort-internal.}
\label{tab:intervals}
\centering
\small
\begin{tabular}{@{}llccc@{}}
\toprule
Cohort & Interval policy & Coverage 50/80/90\% & CRPS & Risk--cov. \\
\midrule
\multirow{3}{*}{Dev ($n{=}15$)}
 & Raw samples & 0.302/0.479/0.558 & 0.1516 & 0.0520 \\
 & Temperature only & 0.625/0.802/0.851 & 0.1524 & 0.0523 \\
 & Propagation + temp. & 0.622/\textbf{0.802}/\textbf{0.853} & \textbf{0.1503} & \textbf{0.0504} \\
\midrule
\multirow{3}{*}{Held-out ($n{=}8$)}
 & Raw samples & 0.281/0.460/0.543 & 0.4016 & 0.0781 \\
 & Temperature only & 0.613/0.813/0.867 & 0.3824 & 0.0781 \\
 & Propagation + temp. & 0.610/\textbf{0.810}/\textbf{0.864} & \textbf{0.3800} & 0.0786 \\
\bottomrule
\end{tabular}
\end{table}

This is also where the diffusion sampler earns its place over deterministic alternatives. In a separate analysis of the same development episodes (with a different aggregation, so absolute values are not comparable across analyses), the diffusion sample distribution beat an ensemble of the deterministic network on CRPS (0.1529 versus 0.1548) and on risk--coverage area (0.0920 versus 0.1129), and after conformal adjustment its 80\% held-out-fold coverage was 0.765 against 0.652 for the ensemble. Combined with Section~\ref{sec:res-paired}, the division of labor is clear: the calibrated operator delivers the point-estimate gain, and the diffusion model turns that calibrated knowledge into the better predictive distribution. Figure~\ref{fig:uq} specifies the interval and representation visualizations.

\subsection{Downstream utility: cleaning that does not cost the task}
\label{sec:res-downstream}

Reconstruction error is not the deployment currency of a BCI; decoding is. On the dataset's native tasks, decoded with deliberately shallow deterministic decoders under a pre-registered grid, the calibrated system passes the test that matters and fails none. On the three-class SSVEP task (CCA, participant-level accuracy; raw 0.773, far from ceiling under ambulation), every method lands within $\pm 0.006$ of raw: matched calibration $-0.001$ ($-0.008$ to $+0.005$) --- ocular cleaning neither helps nor hurts an occipital task, including in the highest EOG-contamination tertile ($-0.007$, $-0.022$ to $+0.007$), where our pre-registered expectation of a gain did not materialize and is reported as such. The ERP oddball is the structured case: the \emph{unguided} pass is the only arm significantly below raw (AUC $-0.0066$, $-0.0107$ to $-0.0030$; worse for 14 of 15 participants) --- cleaning without subject knowledge costs task information --- and every guided variant restores it (matched $+0.003$ against raw; matched over unguided $+0.0096$, $-0.003$ to $+0.020$, better for 12 of 15). The ERP-preservation check completes the picture: matched calibration retains the low-contamination ERP average at 0.931 correlation with raw, on par with ICA (0.933) and far above ASR (0.788) and the eye-subspace projection (0.800). The system does not trade decoding for cleaning; where cleaning does cost task information, the calibrated guide is what pays it back; and the aggressive classical anchors' retention deficit is visible in the ERP itself.

The held-out single pass reproduces the no-harm conclusion on both tasks (SSVEP $-0.002$, interval $-0.012$ to $+0.007$; ERP AUC $-0.002$, interval $-0.009$ to $+0.010$; ERP preservation 0.966), with the guided-versus-unguided direction preserved but not adjudicable at $n=8$. It also sharpens the point that decoding accuracy, waveform preservation, and retention are different currencies: on held-out users ASR's ERP AUC is significantly \emph{above} raw ($+0.042$, interval $+0.011$ to $+0.075$) while its ERP preservation collapses to 0.637 --- aggressive variance removal can help a linear classifier while destroying the morphology it decodes from --- and the eye-subspace projection significantly harms SSVEP decoding ($-0.010$, interval $-0.013$ to $-0.006$). The calibrated system is the one method in the comparison that stays healthy in all three currencies at once, which is the claim a deployable artifact-removal loop actually needs.

\subsection{Calibrated subject information generalizes across interfaces}
\label{sec:res-representation}

The supporting EEGEyeNet experiment asks whether the benefit of matched calibration is tied to the propagation-matrix interface or to the information the calibration carries. Delivered instead through a support-fitted 32-dimensional representation conditioning a deterministic network, participant-matched calibration reduced $\rrmse$ by 0.0608 relative to a zero representation (0.0406--0.0835; 14 of 14 participants) in the plausibility-restricted analysis, and a mismatched representation was harmful, with a matched-minus-mismatched separation of 0.1623 (0.1314--0.1946; 14 of 14). The complete 15-participant analysis was inconclusive ($-$1.428, interval $-$4.415--0.076) because one participant's support episodes carried an artifact-to-signal ratio of 167 against 0.33--4.02 for all others; restricting injected ratios to $[0.1,20]$ gives the analysis above. Since the representation is fitted with paired clean calibration targets, these numbers upper-bound a learned calibration interface rather than certify one. Three conclusions stand. Subject-aware conditioning of a shared model is a sound design across interfaces --- operator or learned representation --- \emph{provided the conditioning is calibrated on the target user}: the benefit is subject-specific (a wrong user's representation harms), which is what a training-time identity code alone cannot deliver for an unseen user. The value of matched calibration is information, not interface: it survives a change of both model class and delivery mechanism. And it is flat in training-population size, changing by $-$0.0224 ($-$0.0491--0.0158) from 30 to 259 training participants: enlarging the population model neither creates nor removes the need to calibrate the individual, which is the premise of this paper measured at scale.

\begin{figure*}[t]
  \centering
  \figureplaceholder{50mm}{\textbf{Intervals and calibration information (design).} Panel A: nominal versus empirical coverage for the three interval policies of Table~\ref{tab:intervals}, development and held-out cohorts side by side with the identity line (held-out arrays in \texttt{t1\_heldout\_uq\_summary.npz}; per-window interval examples for two held-out participants in \texttt{t6\_heldout\_intervals\_sub-01/04.npz}). Panel B: CRPS and risk--coverage area as grouped bars for the three policies per cohort, and, on a separate clearly labeled axis pair, the diffusion-versus-deterministic-ensemble comparison. Panel C: matched-versus-zero representation gain against training-pool size (30, 60, 120, 200, 259) with bootstrap intervals, flat by Section~\ref{sec:res-representation}. Panel D: matched, zero, and mismatched representation contrasts at the largest pool, labeled as the plausibility-restricted analysis with 14 participants. All values from stored aggregates; no invented bands.}
  \caption{Planned figure: interval calibration at no sharpness cost, confirmed on held-out users under frozen settings; the sampler's advantage over a deterministic ensemble; and the flat, interface-independent value of matched calibration.}
  \Description{Placeholder for coverage, CRPS, risk-coverage, and representation-gain panels.}
  \label{fig:uq}
\end{figure*}
